% Generated draft: requires mathematical and source review.
\documentclass[11pt]{article}
\usepackage[margin=25mm]{geometry}
\usepackage{fontspec}
\setmainfont{Latin Modern Roman}
\usepackage{amsmath,amssymb,mathtools,unicode-math}
\setmathfont{Latin Modern Math}
\usepackage{enumitem,needspace,xurl,hyperref}
\setlength{\parindent}{0pt}
\setlength{\parskip}{5pt}
\setlength{\emergencystretch}{3em}
\newcommand{\N}{\mathbb{N}}
\newcommand{\Z}{\mathbb{Z}}
\newcommand{\Q}{\mathbb{Q}}
\newcommand{\R}{\mathbb{R}}
\newcommand{\C}{\mathbb{C}}
\newcommand{\F}{\mathbb{F}}
\newcommand{\A}{\mathbb{A}}
\newcommand{\PP}{\mathbb{P}}
\newcommand{\E}{\mathbb{E}}
\newcommand{\eps}{\varepsilon}
\newcommand{\dd}{\,\mathrm{d}}
\begin{document}
\section*{UnsolvedMath ID: 801189}
\textbf{Source problem code:} catalogue entry
\subsection{Definitions and mathematical statement}
Let \((W,S)\) be a Coxeter system arising from a finite or affine Weyl group, with length function \(\ell:W\to\Z_{\geq 0}\) and Bruhat order \(\leq\). Let \(q\) be an indeterminate. Use the Hecke algebra over \(\Z[q^{1/2},q^{-1/2}]\) with standard basis \(\{T_w:w\in W\}\), the usual braid relations, and quadratic relations
\[
(T_s-q)(T_s+1)=0\qquad(s\in S).
\]
Its bar involution is determined by \(\overline{q^{1/2}}=q^{-1/2}\) and \(\overline{T_w}=T_{w^{-1}}^{-1}\). The Kazhdan--Lusztig polynomials \(P_{u,v}(q)\in\Z[q]\), for \(u\leq v\), are characterized by the bar-invariant elements
\[
C'_v=q^{-\ell(v)/2}\sum_{u\leq v}P_{u,v}(q)T_u,
\qquad \overline{C'_v}=C'_v,
\]
with \(P_{v,v}(q)=1\) and
\[
\deg P_{u,v}\leq\left\lfloor\frac{\ell(v)-\ell(u)-1}{2}\right\rfloor
\qquad(u<v).
\]
Set \(P_{u,v}=0\) when \(u\not\leq v\).

\textbf{Problem (Stanley, supplied attribution, 1999, Problem 16).} Give a direct positive combinatorial interpretation of every coefficient of \(P_{u,v}(q)\), uniformly for finite and affine Weyl groups, with particular emphasis on \(W=S_n\), \(n\geq1\), with its standard adjacent-transposition generators. More explicitly, seek explicitly described finite sets \(\mathcal{C}_{u,v,k}\) of combinatorial objects, for every \(u\leq v\) and integer \(k\geq0\), such that
\[
P_{u,v}(q)=\sum_{k\geq0}|\mathcal{C}_{u,v,k}|q^k.
\]
The description and proof should give a combinatorial explanation of coefficient nonnegativity, without defining the sets merely by their desired cardinalities or presupposing nonnegativity. Only finitely many of these sets may be nonempty for each pair \((u,v)\), and the identity must respect the normalization and degree bound above.

The source does not specify a required class of combinatorial objects or a formal criterion for what counts as a direct interpretation. Its background reports results for special classes, including \(321\)-hexagon-avoiding permutations and Deodhar elements, and describes the general problem as open. These are supplied contextual assessments, not independently verified claims; their exact scope is not specified.
\subsection{Short English statement}
Find explicitly described finite combinatorial objects counted by each coefficient of a Kazhdan--Lusztig polynomial for finite or affine Weyl groups, especially symmetric groups \(S_n\). The counting interpretation should itself explain nonnegativity. The supplied record labels the general problem open and mentions special cases, but those literature assessments require verification.
\subsection{Sources}
\begin{itemize}
\item Supplied local UnsolvedMath registry record: permanent ID \(801189\), exact source code \texttt{catalogue entry}, in the AMR Open Problem Lists collection. The record attributes the problem to Stanley, with proposed year 1999, and labels it \texttt{open}.
\item Stanley, \emph{Positivity Problems and Conjectures in Algebraic Combinatorics} (1999), Problem 16, as cited by the supplied record: \url{https://math.mit.edu/~rstan/papers/problems.pdf}. The supplied provenance records access on 2026-10-16 and extraction method \texttt{pdf+html\_verified}; these are supplied provenance assertions, not a claim of independent access or verification for this draft.
\item The supplied background states that special cases include \(321\)-hexagon-avoiding permutations and Deodhar elements, while the general problem remains open. No separate references or precise theorems supporting these special-case and status assessments are supplied. The registry record was updated on 2026-10-16; its literature assessment is supplied curation and has not been independently verified here.
\end{itemize}
\end{document}
